\chapter{Conclusion}
\label{ch:conclusion}

This thesis set out to assess how Large Language Model (LLM)
support shapes the way novice users produce GDPR Records of
Processing Activities (ROPA). A within-subjects user study of
$73$ participants, each working through all three ropagen
interaction modes (Form, Ask, Chat), was held alongside a
document-quality analysis built from reference-based NLP
metrics and a supplementary LLM-as-a-judge layer. Three
research questions framed the work: how varying degrees of
LLM support shape the novice experience (RQ1), how that
support affects the quality of the resulting documents (RQ2),
and whether a discrepancy exists between user confidence and
actual document quality (RQ3). This chapter summarizes the
answers reached in Chapter~\ref{ch:discussion}, names the
contributions that follow, and closes with what those
contributions imply for the design of LLM-assisted document
tools.


RQ1 asked how varying degrees of LLM support shape the novice
experience during ROPA creation. Less LLM involvement, not
more, produced the better experience. Form was the most
usable, the lightest on cognitive load, the most preferred in
the ranking task, and the fastest to complete; four
independent instruments converged on the same verdict (see
Section~\ref{sec:disc-rq1}). The design intuition the thesis
began with, that adding LLM presence to the interface would
make the task feel easier, does not hold for users who lack
the vocabulary to direct the model. Ask, the hybrid mode,
lost the experience comparison and lost it for a structural
reason: it depends on the user to drive the dialogue, and a
novice without prior GDPR knowledge does not know what to
ask. The mode that asked the user to carry the conversation
failed the user it was designed to help.


RQ2 asked how the redistribution of work affects the quality
of the resulting documents. Each mode plays to a different
strength, and the difference is architectural rather than
incidental. Form supports inclusion: every required field is
present, though some of what fills them was never described
by the user. Chat supports precision: what is on the page
is grounded in the user's scenario, though information may be missing.
Ask trails on both. Two independent
measurement methodologies, reference-based NLP and the
reference-free judge, read the same property from different
sides (see Section~\ref{sec:disc-rq2}). The choice of
interaction mode is therefore also a choice of which kind of
risk a tool prefers to take on: an incomplete document is the
worse legal risk, since it fails the Article~30 obligation
directly; a complete-but-fabricated document is the worse
operational risk, since it commits the organization to
processing details that may not match what it actually does.


RQ3 asked whether what users feel about the AI's output
tracks what the AI actually produces. The answer depends on
which quality measure is laid against perceived confidence.
Users felt clearly less confident in Form than in Ask or
Chat, with the two conversational modes tied at the top of
the perceived confidence composite. Against the
reference-based NLP metrics, that ordering means nothing: the
most and least confident users produce documents whose
BERTScore F1 differs only by rounding. Against the judge's
holistic Overall verdict, the ordering does match. Against
the judge's legal-correctness rubric, the ordering inverts:
the mode users felt most confident about on legal-basis
identification is the mode the judge rated worst on legal
correctness (see Section~\ref{sec:disc-rq3}). Subjective
confidence in a domain the user cannot independently verify
is not a reliable quality signal, and tooling that surfaces
or rewards it without an independent check risks substituting
the appearance of correctness for the fact of it.


The substantive contribution sits in the content-typed
assignment that follows from the chapter-level NLP results.
The thesis is not arguing for a softer synonym of
``hybrid''; it names the criterion. Enumerable sub-records,
where the answer space is a small fixed list such as legal
bases under Article~6(1) GDPR, should be handed to a
checkbox interface. Scenario-specific sub-records, where the
right answer depends on what the organization actually does,
such as retention periods or technical and organizational
measures, should be handed to conversational extraction. The
same principle generalizes beyond ROPA to any structured
legal-document task whose sub-records can be classified along
the enumerable and scenario-specific axes. A separate
architectural result reinforces the point: Legal Correctness
was flat across modes as the generator is the same LLM
in all three, so raising legal-vocabulary quality is a
model-side problem and not a UI-side one. Knowing which dial
moves which dimension is what allows a tool builder to spend
effort where it returns yield.

The methodological contribution is the pairing itself. A
user-experience battery, a reference-based NLP evaluation,
and a reference-free judge layer were held alongside each
other rather than read in isolation, and the
confidence-quality reading reported under RQ3 is the kind of
finding that no single stream could have reached. The
follow-ons that matter most are to build the content-typed
hybrid as a fourth arm in a replication of the user study, to
extend that replication to a more representative
compliance-officer sample, and to pair retrieval-augmented
generation against the GDPR text with a mandatory expert
sign-off at submission for the dimensions interaction
redesign cannot move.

The picture the thesis leaves behind is one of redistribution
rather than replacement. The LLM is useful as a co-author,
not as a teacher: it can take on the parts of legal
documentation the novice cannot supply, but it cannot give
the novice the domain understanding the task presumes.
Neither a novice nor an LLM alone suffices in the making of a
domain-specific document like a ROPA. The user supplies the
account of the processing scenario in their own terms; the
model supplies the vocabulary, scaffolding, and curation the
user cannot reasonably produce; and the document that
satisfies Article~30 exists only where the two meet.
